% Upload this file (or the whole docs/overleaf/ folder) to Overleaf.
% Compile with pdfLaTeX.
\documentclass[11pt]{article}

\usepackage[margin=1in]{geometry}
\usepackage{amsmath,amssymb,mathtools}
\usepackage{booktabs}
\usepackage{colortbl}
\usepackage{xcolor}
\usepackage{tcolorbox}
\usepackage{hyperref}

\hypersetup{colorlinks=true, linkcolor=blue!60!black, citecolor=blue!60!black}

\tcbset{enhanced, boxrule=0.6pt, arc=2pt, left=8pt, right=8pt, top=6pt, bottom=6pt}

\newtcolorbox{samebox}{
  colback=green!6,
  colframe=green!50!black,
  title={\textbf{Same object} (only notation / transpose)},
}

\newtcolorbox{diffbox}{
  colback=orange!8,
  colframe=orange!70!black,
  title={\textbf{Different appearance} (shape or symbol role)},
}

\newcommand{\noise}{\sigma^2}
\newcommand{\Sig}{\Sigma}
\newcommand{\Ph}{\Phi}
\newcommand{\Zmat}{Z}

\title{Woodbury Form for RFF Gaussian Process Marginal Likelihood\\[0.4em]
\large Equivalence of Two $\Sig^{-1}$ Writings (GPPlus notation)}
\author{GPPlus --- \texttt{rff\_mll.py} / \texttt{rff\_utils.py}}
\date{\today}

\begin{document}
\maketitle

\begin{abstract}
This note shows why the Woodbury inverse written with a matrix $\Zmat$ (features $\times$ points)
matches the form in \texttt{gpplus/utils/rff\_utils.py} with $\Ph$ (points $\times$ features).
It also explains \textbf{where the $(\noise I)^{-1}$ factors go} in code (they are never Cholesky-factorized)
and \textbf{why only the $m\times m$ middle matrix $M$} gets a Cholesky decomposition in the solve and MLL.
Green boxes mark \emph{identical} mathematics; orange boxes mark \emph{different} layout or symbols.

A final section summarizes \textbf{Orthogonal Random Features (ORF)} as implemented in GPPlus when
\texttt{RFFKernel(orthogonal=True)}: full Yu et al.\ sampling $\mathbf S\mathbf Q$ with $\chi(d)$ scaling,
independent blocks when $D>d$, and the same cos/sin featurization and Woodbury pipeline as standard RFF.
\end{abstract}

\tableofcontents
\newpage

%=============================================================================
\section{Target covariance (both equations describe the same $\Sig$)}
%=============================================================================

After subtracting the GP mean, the marginal covariance of $y\in\mathbb{R}^n$ is
\begin{equation}
  \label{eq:sigma}
  \boxed{\;
    \Sig \;=\; \noise\, I_n \;+\; \text{(low-rank term)}
  \;}
\end{equation}
with homoskedastic noise variance $\noise$ (GPPlus: \texttt{likelihood.noise}).

For RFF, the low-rank term is a Gram matrix of scaled features.
The fork is only \textbf{matrix layout}: features-as-rows ($\Ph$) vs features-as-columns ($\Zmat$).

%=============================================================================
\section{Two writings of $\Sig^{-1}$ (side by side)}
%=============================================================================

\begin{diffbox}
\textbf{Your equation} ($\Zmat\in\mathbb{R}^{m\times n}$ so the middle factor is $m\times m$):
\begin{equation}
  \label{eq:user}
  \Sig^{-1}
  = (\noise I_n)^{-1}
  - (\noise I_n)^{-1}\,
    \Zmat^{\top}\,
    \underbrace{\Bigl(I_m^{-1} + \Zmat\,(\noise I_n)^{-1}\,\Zmat^{\top}\Bigr)^{-1}}_{M}\,
    \Zmat\,
    (\noise I_n)^{-1}.
\end{equation}
\end{diffbox}

\begin{diffbox}
\textbf{Docstring / \texttt{rff\_utils}} ($\Ph\in\mathbb{R}^{n\times m}$, one row per training point):
\begin{equation}
  \label{eq:code}
  \Sig^{-1}
  = (\noise I_n)^{-1}
  - (\noise I_n)^{-1}\,
    \Ph\,
    \underbrace{\Bigl(I_m + \Ph^{\top}\,(\noise I_n)^{-1}\,\Ph\Bigr)^{-1}}_{M}\,
    \Ph^{\top}\,
    (\noise I_n)^{-1}.
\end{equation}
\end{diffbox}

\begin{samebox}
\textbf{Similarities}
\begin{itemize}
  \item Same $\Sig$: Woodbury inverse of $\noise I_n$ plus a rank-$\le m$ term.
  \item Same middle matrix $M\in\mathbb{R}^{m\times m}$ (Section~\ref{sec:link}).
  \item Same work: Cholesky $M$ in $O(m^3)$, not Cholesky $\Sig$ in $O(n^3)$.
  \item Same MLL log-determinant: $\log|\Sig| = n\log\noise + \log|M|$ --- noise term explicit, $\log|M|$ from Cholesky of $M$ (Section~\ref{sec:logdet}).
\end{itemize}
\end{samebox}

\begin{diffbox}
\textbf{Differences (transpose / naming only)}
\begin{center}
\renewcommand{\arraystretch}{1.25}
\begin{tabular}{@{}lll@{}}
\toprule
 & \textbf{Your $\Zmat$} ($m\times n$) & \textbf{Code $\Ph$} ($n\times m$) \\
\midrule
Layout & features $\times$ points & points $\times$ features \\
Link & $\Ph = \Zmat^{\top}$ & $\Zmat = \Ph^{\top}$ \\
Low-rank part in $\Sig$ & $\Zmat^{\top}\Zmat$ & $\Ph\Ph^{\top}$ \\
Factor after $(\noise I)^{-1}$ (left) & $\Zmat^{\top}$ & $\Ph$ \\
Factor before $(\noise I)^{-1}$ (right) & $\Zmat$ & $\Ph^{\top}$ \\
Middle matrix & $I + \Zmat(\noise I)^{-1}\Zmat^{\top}$ & $I + \Ph^{\top}(\noise I)^{-1}\Ph$ \\
\bottomrule
\end{tabular}
\end{center}
Each orange row is a transpose pair; the last row is the \emph{same} $M$.
\end{diffbox}

%=============================================================================
\section{Why \eqref{eq:user} and \eqref{eq:code} are the same}
%=============================================================================
\label{sec:link}

Define
\begin{equation}
  \Ph \;:=\; \Zmat^{\top} \in \mathbb{R}^{n\times m}.
\end{equation}

\subsection{Covariance}

\begin{samebox}
\[
  \noise I_n + \Ph\Ph^{\top}
  \;=\; \noise I_n + \Zmat^{\top}\Zmat
  \;=\; \Sig.
\]
\end{samebox}

\subsection{Middle matrix $M$}

Your middle:
\[
  I_m + \Zmat\,(\noise I_n)^{-1}\,\Zmat^{\top}
  = I_m + \frac{1}{\noise}\,\Zmat\Zmat^{\top}.
\]
Code middle:
\[
  I_m + \Ph^{\top}\,(\noise I_n)^{-1}\,\Ph
  = I_m + \frac{1}{\noise}\,\Ph^{\top}\Ph.
\]
Since $\Ph^{\top}\Ph = \Zmat\Zmat^{\top}$,
\begin{samebox}
\[
  \boxed{M \;=\; I_m + \frac{\Zmat\Zmat^{\top}}{\noise}
       \;=\; I_m + \frac{\Ph^{\top}\Ph}{\noise}
       \;=\; I_m + \Ph^{\top}(\noise I_n)^{-1}\Ph.}
\]
Note $I_m^{-1}=I_m$, so your $I^{-1}+\cdots$ equals the code's $I+\cdots$.
\end{samebox}

\subsection{Correction term}

Your correction:
\[
  (\noise I_n)^{-1}\,\Zmat^{\top}\,M^{-1}\,\Zmat\,(\noise I_n)^{-1}.
\]
With $\Zmat^{\top}=\Ph$ and $\Zmat=\Ph^{\top}$:
\[
  (\noise I_n)^{-1}\,\Ph\,M^{-1}\,\Ph^{\top}\,(\noise I_n)^{-1},
\]
the second term in \eqref{eq:code}. Hence \eqref{eq:user} $=$ \eqref{eq:code}.

%=============================================================================
\section{Derivation from Woodbury}
%=============================================================================

\begin{theorem}[Woodbury]
For conformable matrices with $A,C$ invertible:
\begin{equation}
  \label{eq:woodbury}
  (A + U C V)^{-1}
  = A^{-1} - A^{-1} U \bigl(C^{-1} + V A^{-1} U\bigr)^{-1} V A^{-1}.
\end{equation}
\end{theorem}

\begin{proof}[Sketch]
Let $M := C^{-1} + V A^{-1} U$. Multiply the RHS of \eqref{eq:woodbury} on the right by $(A+UCV)$:
\[
\begin{aligned}
  &\bigl[A^{-1} - A^{-1} U M^{-1} V A^{-1}\bigr](A + U C V) \\
  &= I + A^{-1}UCV - A^{-1} U M^{-1} V - A^{-1} U (M - C^{-1}) C V \\
  &= I + A^{-1}UCV - A^{-1} U M^{-1} V - A^{-1} U C V + A^{-1} U M^{-1} V = I.
\end{aligned}
\]
\end{proof}

\subsection{Code convention: $\Sig = \noise I_n + \Ph\Ph^{\top}$}

Take $A=\noise I_n$, $U=\Ph$, $C=I_m$, $V=\Ph^{\top}$. Then \eqref{eq:woodbury} is \eqref{eq:code}.

\subsection{Your convention: $\Sig = \noise I_n + \Zmat^{\top}\Zmat$}

With $\Ph=\Zmat^{\top}$, the same choice of $A,U,C,V$ yields \eqref{eq:user}.

%=============================================================================
\section{Where each part of $\Sig^{-1}$ goes}
\label{sec:split}
%=============================================================================

Woodbury does \emph{not} replace $\Sig^{-1}$ by $M^{-1}$ alone. It \textbf{splits} the inverse into a
\textbf{diagonal part} (cheap) and a \textbf{low-rank correction} (compressed into $M$).

\subsection{Term-by-term decomposition (code convention)}

Write \eqref{eq:code} as
\begin{equation}
  \label{eq:split}
  \boxed{
  \Sig^{-1}
  = \underbrace{(\noise I_n)^{-1}}_{\text{(A) diagonal --- no Cholesky}}
  - \underbrace{(\noise I_n)^{-1}\,\Ph}_{\text{(B) matvec}}
  \;\underbrace{M^{-1}}_{\text{(C) only $m\times m$ solve}}
  \;\underbrace{\Ph^{\top}(\noise I_n)^{-1}}_{\text{(D) matvec}}.
  }
\end{equation}

\begin{center}
\renewcommand{\arraystretch}{1.3}
\begin{tabular}{@{}clll@{}}
\toprule
\textbf{Piece} & \textbf{Math} & \textbf{Cost} & \textbf{In GPPlus} \\
\midrule
(A) & $(\noise I)^{-1}$ & $O(n)$ divides & \texttt{b / noise}, \texttt{inv\_noise\_b} \\
(B),(D) & $\Ph$, $\Ph^{\top}$ & $O(nm)$ matvecs & \texttt{z\_train}, \texttt{z\_train.T @ \ldots} \\
(C) & $M^{-1}v$ & $O(m^2)$ after $O(m^3)$ Cholesky & \texttt{cholesky\_solve(..., chol)} \\
\midrule
\multicolumn{4}{@{}l@{}}{%
\textbf{Not used:} $n\times n$ Cholesky or LU of $\Sig=\noise I_n+\Ph\Ph^{\top}$ ($O(n^3)$).} \\
\bottomrule
\end{tabular}
\end{center}

\begin{samebox}
\textbf{The ``other part'' $(\noise I)^{-1}$ is still in the formula.}
It appears as:
\begin{itemize}
  \item the \textbf{leading term} $(\noise I)^{-1}$ (vector: divide by $\noise$);
  \item the \textbf{left and right} factors wrapping $M^{-1}$ in the correction.
\end{itemize}
Because $\noise I_n$ is diagonal, $(\noise I_n)^{-1}$ is \emph{exact division by $\noise$}---no matrix factorization.
\end{samebox}

\begin{diffbox}
\textbf{Common confusion:} ``Woodbury only uses $M$'' means \emph{the expensive matrix inverse / Cholesky
is applied only to $M$} ($m\times m$), not that $(\noise I)^{-1}$ disappears from $\Sig^{-1}$.
\end{diffbox}

\subsection{Structural picture}

\[
  \Sig
  = \underbrace{\noise I_n}_{\text{diagonal (noise)}}
  + \underbrace{\Ph\Ph^{\top}}_{\text{rank $\le m$ (RFF features)}}.
\]

Woodbury inverts this sum by:
\begin{enumerate}
  \item inverting the diagonal piece analytically ($A^{-1}$);
  \item pushing $\Ph\Ph^{\top}$ into the $m\times m$ matrix $M = I + \Ph^{\top}(\noise I)^{-1}\Ph$.
\end{enumerate}

%=============================================================================
\section{Applying $\Sig^{-1}$ to a vector: \eqref{eq:solve}}
\label{sec:solve}
%=============================================================================

For $b\in\mathbb{R}^n$, expanding \eqref{eq:split} acting on $b$ gives
\begin{equation}
  \label{eq:solve}
  \Sig^{-1} b
  = \underbrace{\frac{b}{\noise}}_{\text{(A) } (\noise I)^{-1}b}
  - \underbrace{\frac{1}{\noise}\,\Ph\,M^{-1}\!\left(\frac{\Ph^{\top} b}{\noise}\right)}_{\text{(B)--(C)--(D) correction}}.
\end{equation}

Define the algorithm (matches \texttt{woodbury\_solve\_from\_chol}):
\begin{align}
  \texttt{inv\_noise\_b} &= (\noise I_n)^{-1} b &&= b/\noise, \label{step:a}\\
  \texttt{middle\_rhs} &= \Ph^{\top}\,(\noise I_n)^{-1} b &&= \Ph^{\top}\,\texttt{inv\_noise\_b}, \label{step:bd}\\
  \texttt{inner} &= M^{-1}\,\texttt{middle\_rhs} &&\text{Cholesky on $M$ only}, \label{step:c}\\
  \texttt{correction} &= (\noise I_n)^{-1}\,\Ph\,\texttt{inner} &&= \Ph\,\texttt{inner}/\noise, \label{step:b2}\\
  \Sig^{-1}b &= \texttt{inv\_noise\_b} - \texttt{correction}. \label{step:final}
\end{align}

\begin{samebox}
\textbf{Only step \eqref{step:c} uses the Cholesky factor} \texttt{chol} from \texttt{woodbury\_factor}
($LL^{\top}=M$). Steps \eqref{step:a}--\eqref{step:bd} and \eqref{step:b2} implement the
$(\noise I)^{-1}$ and $\Ph,\Ph^{\top}$ factors explicitly.
\end{samebox}

\texttt{woodbury\_factor}: \texttt{cholesky(}$M$\texttt{)} --- not \texttt{cholesky(}$\Sig$\texttt{)}.

\begin{diffbox}
Avoid: ``Cholesky Woodbury on $\Sig$'' (sounds $n\times n$).
\end{diffbox}
\begin{samebox}
Use: ``Woodbury for $\Sig^{-1}$: divide by $\noise$ for $(\noise I)^{-1}$, Cholesky $M$ for the low-rank correction''.
\end{samebox}

%=============================================================================
\section{Marginal likelihood: why MLL also splits noise vs.\ $M$}
\label{sec:logdet}
%=============================================================================

Marginal log-density for centered $y$:
\begin{equation}
  \label{eq:mll}
  \log p(y)
  = -\tfrac{1}{2}\,\underbrace{y^{\top}\Sig^{-1}y}_{\text{quadratic --- Section~\ref{sec:solve}}}
  -\tfrac{1}{2}\,\underbrace{\log|\Sig|}_{\text{log-det --- below}}
  -\tfrac{n}{2}\log(2\pi).
\end{equation}

\subsection{Quadratic term $y^{\top}\Sig^{-1}y$}

Compute $\alpha=\Sig^{-1}y$ with \eqref{eq:solve}--\eqref{step:final}: one Cholesky of $M$, plus
$O(n)$ divides and $O(nm)$ matvecs for the $(\noise I)^{-1}$ and $\Ph$ factors.
Then $y^{\top}\Sig^{-1}y = y^{\top}\alpha$ (\texttt{woodbury\_quadratic\_form}).

\subsection{Log determinant $\log|\Sig|$}

\begin{lemma}[Matrix determinant lemma]
For $\Ph\in\mathbb{R}^{n\times m}$,
\[
  |\Sig|
  = |\noise I_n + \Ph\Ph^{\top}|
  = (\noise)^n \,\bigl|I_m + \Ph^{\top}\Ph/\noise\bigr|
  = \noise^n\,|M|.
\]
\end{lemma}

So $\log|\Sig|$ \textbf{splits the same way} as $\Sig^{-1}$:
\begin{equation}
  \label{eq:logdet-split}
  \boxed{
  \log|\Sig|
  = \underbrace{n\log\noise}_{\text{diagonal / noise part --- scalar, no Cholesky}}
  + \underbrace{\log|M|}_{\text{from Cholesky of $M$ only}}.
  }
\end{equation}

In code (\texttt{woodbury\_log\_det\_from\_chol}):
\begin{itemize}
  \item \textbf{Noise part:} \texttt{n * noise.log()} implements $n\log\noise$.
  \item \textbf{$M$ part:} \texttt{2 * sum(log diag(chol))} implements $\log|M|$ from $LL^{\top}=M$.
\end{itemize}

\texttt{woodbury\_marginal\_log\_likelihood} uses \emph{one} Cholesky of $M$ for both
$y^{\top}\Sig^{-1}y$ and $\log|M|$; the $(\noise I)^{-1}$ contribution to the MLL is
\textbf{division} (in the solve) and \textbf{$n\log\noise$} (in the log-det), not an $n\times n$ factorization.

\begin{samebox}
\textbf{Summary for MLL}
\begin{center}
\renewcommand{\arraystretch}{1.25}
\begin{tabular}{@{}lll@{}}
\toprule
\textbf{MLL ingredient} & \textbf{Uses $(\noise I)^{-1}$?} & \textbf{Cholesky?} \\
\midrule
$y^{\top}\Sig^{-1}y$ & yes (divide by $\noise$ in \eqref{step:a}--\eqref{step:b2}) & only $M$ ($m\times m$) \\
$\log|\Sig|$ & yes ($n\log\noise$ term) & only $\log|M|$ from \texttt{chol} \\
\bottomrule
\end{tabular}
\end{center}
\end{samebox}

%=============================================================================
\section{Orthogonal Random Features (ORF)}
\label{sec:orf}
%=============================================================================

Woodbury (Sections~\ref{sec:split}--\ref{sec:logdet}) applies to \emph{any} feature matrix $\Ph\in\mathbb{R}^{n\times m}$
with $m=2D$. This section describes how GPPlus \emph{draws} the underlying RFF frequencies when
\texttt{orthogonal=True} on \texttt{RFFKernel} / \texttt{RFFGPR}.
Implementation: \texttt{\_sample\_orf\_weights} and \texttt{init\_rbf\_weights} in
\texttt{gpplus/utils/rff\_utils.py}.
Reference: Yu, Suresh, Choromanski et al., ``Orthogonal Random Features'',
\href{https://arxiv.org/abs/1610.09072}{arXiv:1610.09072}.

\subsection{Standard RFF (i.i.d.\ frequencies)}

For input dimension $d$ and $D$ frequencies, GPPlus stores a weight matrix
$W\in\mathbb{R}^{d\times D}$ whose \textbf{columns} are $\mathbf w_j\in\mathbb{R}^d$.
Standard RFF draws
\begin{equation}
  \label{eq:rff-iid}
  \mathbf w_j \stackrel{\text{i.i.d.}}{\sim} \mathcal N(0, I_d),
  \qquad j=1,\ldots,D.
\end{equation}
Scaled features (before outputscale in \texttt{RFFGPR}) are
\begin{equation}
  \label{eq:featurize}
  \mathbf z(x)
  = \frac{1}{\sqrt{D}}
  \begin{bmatrix}
    \cos(\tilde{\mathbf w}_1^{\top}\tilde{\mathbf x}) \\
    \vdots \\
    \cos(\tilde{\mathbf w}_D^{\top}\tilde{\mathbf x}) \\
    \sin(\tilde{\mathbf w}_1^{\top}\tilde{\mathbf x}) \\
    \vdots \\
    \sin(\tilde{\mathbf w}_D^{\top}\tilde{\mathbf x})
  \end{bmatrix}
  \in \mathbb{R}^{2D},
\end{equation}
where $\tilde{\mathbf x}$ applies GPPlus lengthscale scaling ($10^{\ell/2}$ per dimension in
\texttt{featurize\_rbf}) and $\tilde{\mathbf w}_j$ may include optional ARD scaling on $W$
when \texttt{resample\_weights(spectral=True)}.
One training row is $\Ph_i = \mathbf z(x_i)^{\top}$ (with outputscale applied in \texttt{RFFGPR}).

\subsection{Full ORF (Yu et al., Eq.~(2))}

Plain QR on a Gaussian matrix gives \emph{unit-norm} directions; Gaussian RFF rows have
$\chi(d)$-distributed norms. Full ORF restores that distribution while keeping directions orthogonal
within each block.

\begin{equation}
  \label{eq:orf-paper}
  \boxed{
  W_{\mathrm{ORF}} = \frac{1}{\sigma}\, S\, Q,
  \qquad
  \mathbf w_j = s_j\,\mathbf q_j,
  }
\end{equation}
where:
\begin{itemize}
  \item $Q\in\mathbb{R}^{d\times d}$ is a random orthogonal matrix from the QR decomposition of a
        $d\times d$ Gaussian matrix (uniform on the Stiefel manifold);
  \item $S=\mathrm{diag}(s_1,\ldots,s_D)$ with $s_j \sim \chi(d)$ i.i.d.;
  \item $\mathbf q_j$ is the $j$-th \textbf{column} of $Q$ (orthonormal directions in $\mathbb{R}^d$).
\end{itemize}
Each $\mathbf w_j$ is marginally $\mathcal N(0,I_d)$, so kernel approximation remains unbiased;
the variance of $\hat K(x,y)$ is lower than i.i.d.\ RFF for the same $D$ (Theorem~1 in the paper).

\begin{diffbox}
\textbf{ORF$'$ (simplified, not implemented):}
$W=\frac{\sqrt{d}}{\sigma}\,Q$ without $\chi(d)$ scaling uses unit-norm rows/columns only.
GPPlus implements \textbf{full ORF} ($S$ included), not ORF$'$.
\end{diffbox}

\subsection{When $D > d$: independent blocks}

When more than $d$ frequencies are needed, the paper repeats Eq.~\eqref{eq:orf-paper}
\textbf{independently} for each block of size $d$ (Corollary~1).
GPPlus concatenates columns:
\begin{equation}
  W = \big[\, W^{(1)} \;\|\; W^{(2)} \;\|\; \cdots \,\big],
  \qquad W^{(b)} \in \mathbb{R}^{d\times d},
\end{equation}
each $W^{(b)}$ from a fresh Gaussian $\to$ QR $\to$ $\chi(d)$ scales.

\subsection{GPPlus mapping (code convention)}

\begin{center}
\renewcommand{\arraystretch}{1.25}
\begin{tabular}{@{}ll@{}}
\toprule
\textbf{Concept} & \textbf{GPPlus} \\
\midrule
Input dimension & $d = \texttt{num\_dims}$ \\
Number of frequencies & $D = \texttt{num\_samples} = \texttt{num\_rff}$ \\
Feature dimension & $m = 2D$ (cos/sin pairs unchanged) \\
Weight buffer & \texttt{randn\_weights} $\in \mathbb{R}^{d\times D}$, column $j$ is $\mathbf w_j$ \\
ORF enabled & \texttt{RFFKernel(orthogonal=True)}, \texttt{RFFGPR(..., orthogonal=True)} \\
Sampling & \texttt{\_sample\_orf\_weights} $\to$ \texttt{init\_rbf\_weights(..., orthogonal=True)} \\
$\chi(d)$ draw & $\|\mathcal N(0,I_d)\|$ (norm of a $d$-dim.\ standard Gaussian) \\
ARD on weights & Applied \textbf{after} ORF: $W \leftarrow W \odot (1/\boldsymbol{\ell})$ if \texttt{lengthscale} given \\
Kernel bandwidth on $x$ & Still in \texttt{featurize\_rbf} via $10^{\ell/2}$; paper's $1/\sigma$ is not duplicated \\
Multi-init draws & \texttt{RFFParameterInitializer} pre-draws one $W$ per init with the same ORF rule \\
\bottomrule
\end{tabular}
\end{center}

Algorithm (one block of size $b\le d$):
\begin{align}
  G &\sim \mathcal N(0,I)^{d\times d}, &
  Q &\leftarrow \mathrm{qr}(G), \\
  s_j &\sim \chi(d), &
  \mathbf w_j &\leftarrow s_j\,\mathbf q_j, \quad j=1,\ldots,b.
\end{align}

\begin{samebox}
\textbf{What ORF does \emph{not} change}
\begin{itemize}
  \item Featurization \eqref{eq:featurize}: still cos/sin with $1/\sqrt{D}$ scaling.
  \item Woodbury MLL, $\Sig$, $M$, and all formulas in Sections~\ref{sec:split}--\ref{sec:logdet}.
  \item Only the \textbf{random draw} of $W$ differs from \eqref{eq:rff-iid}.
\end{itemize}
\end{samebox}

\begin{diffbox}
\textbf{Not implemented (out of scope):} SORF (structured orthogonal features via Walsh--Hadamard
matrices, Eq.~(5) in the paper) and explicit materialization-free $O(d\log d)$ transforms.
GPPlus always stores explicit $W\in\mathbb{R}^{d\times D}$.
\end{diffbox}

%=============================================================================
\section{Highlight map (one page)}
%=============================================================================

\begin{center}
\renewcommand{\arraystretch}{1.2}
\begin{tabular}{@{}p{0.44\textwidth}p{0.44\textwidth}@{}}
\toprule
\multicolumn{1}{c}{\cellcolor{orange!15}\textbf{Different look}} &
\multicolumn{1}{c}{\cellcolor{green!12}\textbf{Same math}} \\
\midrule
Symbol $\Zmat$ vs $\Ph$ & $\Ph=\Zmat^{\top}$ \\
Order: $\Zmat^{\top} M^{-1}\Zmat$ vs $\Ph M^{-1}\Ph^{\top}$ & same bilinear form \\
$I^{-1}+\cdots$ vs $I+\cdots$ & both $= M$ \\
\midrule
\multicolumn{2}{@{}c@{}}{\textbf{Shared}} \\
\midrule
\multicolumn{2}{@{}l@{}}{$\Sig = \noise I_n + \text{rank-}\le m$ term} \\
\multicolumn{2}{@{}l@{}}{$M$ is $m\times m$; factor $M$, not $\Sig$} \\
\multicolumn{2}{@{}l@{}}{$\Sig^{-1}b$: $(\noise I)^{-1}$ by divide; $M^{-1}$ by Cholesky} \\
\multicolumn{2}{@{}l@{}}{$\log|\Sig|$: $n\log\noise$ explicit $+$ $\log|M|$ from Cholesky} \\
\midrule
\multicolumn{2}{@{}l@{}}{ORF: $W=SQ$, $\chi(d)$ scales; same $\Ph$ Woodbury (Section~\ref{sec:orf})} \\
\bottomrule
\end{tabular}
\end{center}

\vspace{1em}
\noindent\texttt{gpplus/utils/rff\_utils.py}, \texttt{gpplus/training/rff\_mll.py}.

\end{document}
